\chapter{Marco contextual}\label{cap:contexto}

Este capítulo sitúa el problema técnico de la tesis en su contexto: qué es el lavado de dinero
y cómo se organiza (Sección \ref{sec:ctx_lavado}), cómo lo transforman las criptomonedas
(Sección \ref{sec:ctx_cripto}), cómo lo detectan hoy las instituciones financieras y por qué
esos sistemas son insuficientes (Sección \ref{sec:ctx_monitoreo}), y qué exige la regulación a
los sistemas basados en inteligencia artificial (Sección \ref{sec:ctx_regulacion}). Cierra con
la descripción del conjunto de datos Elliptic como representación de ese problema (Sección
\ref{sec:ctx_elliptic}).

\section{El lavado de dinero}\label{sec:ctx_lavado}

El lavado de dinero es el conjunto de operaciones con las que se oculta el origen ilícito de
unos fondos para que puedan usarse como si fueran legítimos. Es el paso que permite a una
actividad criminal convertir sus ganancias en patrimonio aprovechable, y por eso combatirlo es
una forma de combatir los delitos que lo originan. Es un problema global desde hace décadas y una amenaza para la economía y la sociedad \cite{Kute2021AMLReview}.

\subsection{Las tres etapas}

La descripción más extendida del proceso distingue tres etapas \cite{Irwin2012Typologies}:

\begin{enumerate}
  \item \textbf{Colocación.} Los fondos de origen ilícito entran por primera vez al sistema
        financiero, por ejemplo mediante depósitos fraccionados en efectivo, compras de
        instrumentos financieros o, en el caso de las criptomonedas, la compra de activos
        virtuales en plataformas de intercambio.
  \item \textbf{Estratificación.} Los fondos se mueven a través de muchas operaciones, cuentas
        e intermediarios con el fin de separar el dinero de su origen y dificultar su rastreo.
        Es la etapa en la que se generan las cadenas de transacciones que un análisis de grafos
        puede seguir.
  \item \textbf{Integración.} Los fondos vuelven a la economía como ingresos aparentemente
        legítimos, por ejemplo a través de inversiones, compras de bienes o empresas.
\end{enumerate}

Irwin et al. analizaron 184 tipologías documentadas por organismos de prevención del lavado y
del financiamiento del terrorismo, y encontraron que quienes lavan dinero y quienes financian
el terrorismo prefieren técnicas algo distintas en cada etapa, y que cuantas más técnicas se
combinan, más dinero se logra ocultar \cite{Irwin2012Typologies}. Esa combinación de técnicas
es la que vuelve el problema difícil para un sistema de reglas: el patrón no está en una
operación aislada sino en la secuencia y en la red de operaciones.

\subsection{La escala del problema}

Medir el lavado es difícil precisamente porque está diseñado para no ser visto. Las
estimaciones disponibles para las criptomonedas dan una idea de su magnitud. Foley et al.
estimaron, con técnicas de análisis de la cadena de bloques, que cerca de una cuarta parte de
los usuarios de Bitcoin participaba en actividades ilegales, y que esas actividades movían unos
76.000 millones de dólares al año, el 46~\% de las transacciones de Bitcoin, una cifra comparable
a los mercados de drogas ilegales de Estados Unidos y Europa \cite{Foley2019Bitcoin}. Los
mismos autores observaron que la proporción ilegal disminuía a medida que crecía el interés
general en Bitcoin y aparecían criptomonedas más opacas. En otro frente, Paquet-Clouston et al.
reconstruyeron los pagos de 35 familias de \textit{ransomware} en Bitcoin y estimaron un mercado
de al menos 12,7 millones de dólares entre 2013 y mediados de 2017, concentrado en unos pocos
actores \cite{PaquetClouston2019Ransomware}.

\section{Criptomonedas y lavado de dinero}\label{sec:ctx_cripto}

\subsection{Qué cambia con las criptomonedas}

Las criptomonedas permiten transferir valor de persona a persona sin un intermediario
financiero, y la cadena de bloques asocia las transacciones a direcciones y no a identidades.
Albrecht et al. señalan que esas dos propiedades las volvieron atractivas para los mercados de
la red oscura y otras redes criminales, y analizan cómo se incorporan a cada etapa del lavado
\cite{Albrecht2019Crypto}. Almeida et al. analizan los métodos que se usan con criptomonedas, desde los intercambios entre particulares hasta los servicios de mezcla de fondos, en cada una de las tres etapas, y encuentran que las técnicas más sofisticadas dan más anonimato pero exigen más conocimiento técnico \cite{Almeida2025Techniques}. Subashi destaca que buena parte de estas transacciones ocurre fuera del sistema bancario convencional, lo que las hace difíciles de rastrear, y que la respuesta exige coordinación internacional \cite{Subashi2024Crypto}.

El punto débil de quien lava con criptomonedas es la conversión a dinero tradicional, porque
ahí reaparece un intermediario regulado. Dupuis y Gleason identificaron seis «puertas abiertas»
para convertir criptomonedas en moneda tradicional que seguían disponibles para quien busca
lavar fondos, y las analizaron con el paradigma de la dialéctica regulatoria: cada regulación
cierra una puerta y los infractores buscan la siguiente \cite{Dupuis2020OpenDoors}.

\subsection{Seudonimato y trazabilidad}

La cadena de bloques de Bitcoin es pública: cualquiera puede ver todas las transacciones,
sus montos y las direcciones que intervienen. Lo que no es público es quién controla cada
dirección. Esta combinación de transparencia y seudonimato es la que hace posible el análisis
forense de la cadena de bloques y, a la vez, la que lo limita \cite{Pocher2023Anomalous}. Un
investigador puede seguir los fondos de transacción en transacción, pero necesita información
externa, como la de una plataforma de intercambio, para atribuir una dirección a una persona. Por
eso el problema se plantea de forma natural como el análisis de un grafo de transacciones, en
el que lo que se observa son flujos y no identidades \cite{Weber2019Elliptic}.

\section{Sistemas de monitoreo de transacciones}\label{sec:ctx_monitoreo}

\subsection{Monitoreo basado en reglas}

Las instituciones financieras están obligadas a monitorear las transacciones de sus clientes y
reportar las sospechosas. La mayoría lo hace con sistemas de reglas: umbrales de monto,
frecuencias inusuales, contrapartes en listas de riesgo o combinaciones de estos criterios. El
problema conocido de estos sistemas es la proporción de falsas alarmas. Ketenci et al.
reportan que la mayoría de los sistemas en uso son basados en reglas y que más del 90~\% de sus
alertas son falsos positivos \cite{Ketenci2021TimeFrequency}. Oztas et al. entrevistaron a
especialistas en prevención del lavado, que identificaron esa tasa de falsos positivos, con sus
costos e ineficiencias, como uno de los problemas centrales del monitoreo actual, y señalaron el
aprendizaje automático, el análisis de grafos y la detección de anomalías como vías para
superarlo \cite{Oztas2024Monitoring}.

\subsection{Monitoreo con aprendizaje automático}

Los sistemas de monitoreo basados en aprendizaje automático han complementado a los de reglas para reducir las falsas alarmas y el esfuerzo de revisar manualmente todas las alertas \cite{Labanca2022Amaretto}. Las revisiones sobre detección de fraude financiero muestran que el aprendizaje supervisado ha sido el más estudiado, y que sus limitaciones han llevado a explorar enfoques semisupervisados y no supervisados \cite{Hilal2022Anomaly}. Los modelos supervisados suelen detectar mejor, pero necesitan muchas etiquetas; los no supervisados detectan patrones anómalos nuevos, pero producen más falsas alarmas. Labanca et al. combinan ambos en un esquema de aprendizaje
activo en el que el modelo elige qué transacciones conviene que el analista revise, para
aprovechar mejor su retroalimentación \cite{Labanca2022Amaretto}. En todos los casos la
escasez de ejemplos positivos es una constante: los casos confirmados de lavado son muy pocos
frente al volumen de operaciones legítimas, lo que convierte el desbalance de clases en una
característica estructural del problema y no en un accidente de los datos.

Los entrevistados por Oztas et al. incluyeron la explicabilidad entre los requisitos para que
un sistema de monitoreo basado en aprendizaje automático se adopte en la práctica
\cite{Oztas2024Monitoring}. El analista que recibe una alerta tiene que decidir si la reporta, y
para hacerlo necesita saber qué la motivó.

\section{Regulación e inteligencia artificial explicable}\label{sec:ctx_regulacion}

El uso de modelos opacos en decisiones con consecuencias para las personas plantea exigencias
de transparencia que la regulación ha empezado a recoger. Bertrand et al. examinan si los
sistemas de prevención del lavado basados en inteligencia artificial respetan los derechos
fundamentales europeos, y recuerdan que la normativa de la Unión Europea obliga a las entidades
financieras a un monitoreo continuo capaz de detectar transacciones inusuales o sospechosas; la
introducción de algoritmos de aprendizaje automático más opacos en ese monitoreo plantea riesgos
de protección de datos que solo se han estudiado en parte \cite{Bertrand2021Rights}.

En el sector financiero en general, Giudici et al. proponen un marco de indicadores para
medir los riesgos de la inteligencia artificial, en el que traducen los requisitos regulatorios
propuestos para la inteligencia artificial en cuatro principios medibles: sostenibilidad,
exactitud, equidad y explicabilidad \cite{Giudici2023Risk}. Una revisión sistemática de 35
casos de uso de explicabilidad en finanzas encontró que la detección de fraude es una de las
aplicaciones más frecuentes, que SHAP y LIME son los métodos más usados, y que los beneficios
buscados son la transparencia, la confianza de los interesados y el cumplimiento regulatorio
\cite{FinXAI2026Mapping}.

Esta tesis parte de una observación que esa literatura suele dar por supuesta: para que una
explicación sirva como evidencia ante un auditor, debe ser reproducible. Si el mismo método, aplicado a la misma transacción y al mismo modelo, produce explicaciones distintas, difícilmente puede sostener una decisión regulatoria; Gawantka et al. plantean el mismo problema para los métodos de explicabilidad no deterministas \cite{Gawantka2024Explanations}. La estabilidad de las explicaciones es,
por eso, una condición previa a su uso en un entorno regulado.

\section{El conjunto de datos Elliptic}\label{sec:ctx_elliptic}

Los datos reales de lavado de dinero rara vez son públicos, porque contienen información
protegida de clientes y de investigaciones. Elliptic es una excepción: es un grafo de
transacciones de Bitcoin publicado por la empresa del mismo nombre junto con investigadores del
MIT e IBM, con una parte de las transacciones etiquetadas como lícitas o ilícitas según su
asociación con entidades conocidas \cite{Weber2019Elliptic}. Las ilícitas corresponden a
categorías como estafas, \textit{malware}, organizaciones terroristas, \textit{ransomware} y
esquemas Ponzi; las lícitas, a plataformas de intercambio, proveedores de billeteras, mineros y
servicios legales.

Elliptic se ha convertido en la referencia pública para la detección de transacciones
ilícitas en criptomonedas \cite{Alarab2023Graph, DTA2022Evolving}. Tiene tres rasgos que lo
hacen adecuado para esta tesis: es real, está desbalanceado de forma natural y su estructura de
grafo permite aplicar GNN. Tiene también dos limitaciones que condicionan el diseño. La primera es
que sus variables están anonimizadas, de modo que no se puede juzgar si una explicación señala
algo con sentido para un experto; solo se puede medir si es estable. La segunda es el cierre de
un mercado ilegal en el paso de tiempo 43, que cambia la distribución de las transacciones
ilícitas posteriores y que ningún modelo logra capturar con las variables originales
\cite{Weber2019Elliptic, Alarab2023Recurrent}. El Capítulo \ref{cap:metodologia} describe el
conjunto de datos en detalle y la forma en que se trató esa segunda limitación.
